%========================================================
% SECTION 5.1 — Deterministic Phase Sweep
%========================================================

\section{Deterministic Phase Sweep}

%--------------------------------------------------------
% Subsection 5.1.1 — Objective and Physical Motivation
%--------------------------------------------------------

\subsection{Objective and Physical Motivation}

This experiment establishes the baseline propagation regime in which the optical fiber acts as an effective deterministic phase shifter preserving both the purity and the entanglement of the bipartite state.

The objective is to characterize how observable correlations and state-characterization quantities evolve under the reduced unitary fiber model introduced in Section~\ref{sec:reduced_phase_model}.

This experiment represents the minimal physically relevant propagation scenario. The fiber is modeled exclusively through a deterministic differential phase shift acting on the polarization components of the transmitted photons, while all non-unitary effects are neglected.

Consequently, the experiment isolates the role of coherent phase evolution and establishes the baseline regime against which later noisy and depolarizing models will be compared.

%--------------------------------------------------------
% Subsection 5.1.2 — Model and Analytical Predictions
%--------------------------------------------------------

\subsection{Model and Analytical Predictions}

The initial bipartite state is chosen as the Bell state

\[
|\Psi^{+}\rangle
=
\frac{1}{\sqrt{2}}
\left(
|01\rangle + |10\rangle
\right).
\]

Under the reduced phase model, the fiber induces an effective relative phase shift between the polarization components, leading to the evolved state

\begin{equation}
|\psi(\theta)\rangle
=
\frac{1}{\sqrt{2}}
\left(
|01\rangle + e^{i\theta}|10\rangle
\right).
\label{eq:phase_evolved_bell_state}
\end{equation}

The diagonal correlation observables considered in the simulation are

\[
c_{xx}
=
\langle \sigma_x \otimes \sigma_x \rangle,
\qquad
c_{yy}
=
\langle \sigma_y \otimes \sigma_y \rangle,
\qquad
c_{zz}
=
\langle \sigma_z \otimes \sigma_z \rangle.
\]

Direct analytical evaluation gives

\[
c_{xx} = \cos\theta,
\qquad
c_{yy} = \cos\theta,
\qquad
c_{zz} = -1.
\]

Since the evolution is generated by local unitary transformations, the global state remains pure:

\[
\gamma(\rho(\theta))
=
\operatorname{Tr}
\left(
\rho(\theta)^2
\right)
=
1.
\]

The reduced density operators remain maximally mixed:

\[
\rho_A(\theta)
=
\rho_B(\theta)
=
\frac{I}{2},
\]

which implies

\[
\gamma(\rho_A)
=
\gamma(\rho_B)
=
\frac{1}{2}.
\]

The fidelity with respect to the reference Bell state \( |\Psi^+\rangle \) is

\[
F_{\Psi^+}(\rho(\theta))
=
\frac{1+\cos\theta}{2}.
\]

The concurrence remains constant:

\[
C(\rho(\theta))
=
1,
\]

demonstrating the invariance of bipartite entanglement under local unitary evolution.

\medskip

The experiment therefore provides a controlled setting in which observable correlations evolve continuously with the phase parameter \( \theta \), while the underlying entanglement structure remains unchanged.

%--------------------------------------------------------
% Subsection 5.1.3 — Numerical Results
%--------------------------------------------------------

\subsection{Numerical Results}

A numerical sweep over the interval

\[
\theta \in [0,2\pi]
\]

is performed in order to evaluate both correlation observables and state-characterization quantities throughout the phase evolution.

Figure~\ref{fig:experiment1_validation} compares the numerical results obtained from the simulator with the corresponding analytical predictions for the diagonal correlation observables.

\begin{figure}[htbp]
    \centering
    \includegraphics[width=\linewidth]{figures/experiment_1/experiment_1_correlations_validation.png}
    \caption{
    Numerical and analytical comparison for the correlation observables \(c_{xx}\), \(c_{yy}\), and \(c_{zz}\).
    The agreement validates the implementation of the reduced phase model.
    }
    \label{fig:experiment1_validation}
\end{figure}

The numerical curves reproduce the expected cosine dependence for \(c_{xx}\) and \(c_{yy}\), while \(c_{zz}\) remains constant at \(-1\) over the entire propagation range.

The numerical and analytical curves are visually indistinguishable within numerical precision.

\medskip

Figure~\ref{fig:experiment1_state_metrics} shows the evolution of the state-characterization quantities.

\begin{figure}[htbp]
    \centering
    \includegraphics[width=\linewidth]{figures/experiment_1/experiment_1_state_metrics.png}
    \caption{
    State-characterization quantities for the phase-evolved state.
    Global purity and concurrence remain constant, while fidelity varies with the relative phase \( \theta \).
    }
    \label{fig:experiment1_state_metrics}
\end{figure}

The global purity remains equal to unity throughout the evolution, confirming that the propagated state remains pure. The reduced purities remain fixed at \(1/2\), while the concurrence remains equal to \(1\), indicating preservation of maximal entanglement.

In contrast, the fidelity with respect to the reference Bell state varies continuously with the phase parameter.

\medskip

The full correlation tensor obtained numerically is shown in Figure~\ref{fig:experiment1_tensor}.

\begin{figure}[htbp]
    \centering
    \includegraphics[width=\linewidth]{figures/experiment_1/experiment_1_correlation_tensor.png}
    \caption{
    Full correlation tensor \(T_{ij}(\theta)\) for the phase-evolved state.
    The off-diagonal tensor components reveal the rotation of the correlation structure in the transverse plane.
    }
    \label{fig:experiment1_tensor}
\end{figure}

The numerical results confirm the analytical tensor structure

\begin{equation}
T(\theta)
=
\begin{pmatrix}
\cos\theta & -\sin\theta & 0 \\
\sin\theta & \cos\theta & 0 \\
0 & 0 & -1
\end{pmatrix}.
\label{eq:phase_tensor_experiment}
\end{equation}

The non-zero off-diagonal components

\[
T_{xy}
=
-\sin\theta,
\qquad
T_{yx}
=
\sin\theta
\]

demonstrate that the phase evolution induces a rotation of the correlation structure within the transverse \(x\)--\(y\) plane.

This structure is not visible when considering only diagonal observables, but emerges naturally within the full tensor representation.

%--------------------------------------------------------
% Subsection 5.1.4 — Discussion
%--------------------------------------------------------

\subsection{Discussion}

The numerical results are in full agreement with the analytical predictions.

The experiment confirms that deterministic phase evolution modifies measurable two-photon correlations while preserving both the purity and the entanglement of the global state.

In particular, the transverse correlations \(c_{xx}\) and \(c_{yy}\) exhibit the expected cosine dependence on the phase parameter \( \theta \), whereas \(c_{zz}\) remains invariant.

At the same time, the concurrence remains constant despite the continuous variation of the fidelity with respect to the reference Bell state.

This distinction is physically significant.

The fidelity quantifies the overlap with a specific reference state and therefore depends explicitly on the relative phase. The concurrence, instead, measures the intrinsic entanglement content of the bipartite state and remains invariant under local unitary transformations.

Consequently, changes in fidelity do not necessarily imply degradation of entanglement.

The experiment therefore establishes a clear separation between:

\begin{itemize}
    \item modifications of observable correlations induced by coherent phase evolution;
    \item genuine loss of entanglement caused by non-unitary processes.
\end{itemize}

This distinction will become essential in the interpretation of the later experiments involving statistical averaging and depolarization.

%--------------------------------------------------------
% Subsection 5.1.5 — Physical Interpretation
%--------------------------------------------------------

\subsection{Physical Interpretation}

The reduced phase model admits a direct physical interpretation in terms of fiber-induced birefringence.

In optical fibers, the horizontal and vertical polarization components generally propagate with different effective refractive indices. As a consequence, the two polarization modes accumulate different propagation phases.

After compensation of global polarization rotations, the dominant residual effect is a differential phase shift between the polarization components. Within the qubit representation, this residual evolution corresponds to a local unitary rotation around the \(Z\)-axis of the Bloch sphere.

The phase parameter \( \theta \) therefore represents the effective relative phase accumulated during propagation.

The experiment shows that this mechanism modifies measurable correlations without degrading the entanglement of the transmitted state. The global quantum state remains pure, and the bipartite entanglement remains maximal throughout the evolution.

At the same time, the dependence of the transverse correlations on \( \theta \) demonstrates that experimentally accessible observables remain highly sensitive to phase fluctuations.

This experiment therefore establishes the minimal effective model of coherent fiber propagation: a deterministic unitary transformation that continuously rotates the observable correlation structure while preserving the underlying entangled nature of the bipartite quantum state.